\subsection{Asignación fija por recurso y quincena}
Synaptix propone mantener los ocho titulares de SD1 y agregar puestos especializados de ejecución según la Tabla~\ref{tab:sd7-refuerzo-fijo}. El registro asociado \path{asignacion_fija_recursos.csv} identifica paquete, mes, quincena, perfil, recurso y HH. Asigna 20.820 HH de 797 paquetes con mes fijo y conserva aparte 328 HH de 16 paquetes del piloto estacional IN2. El total fijo estimado sigue siendo 21.148 HH: 13.588 de implementación y 7.560 de operación. Los trabajos dependientes de cantidades y la cobertura externa se agregan mediante sus fórmulas; no están incluidos en esta asignación.

Los perfiles son A: Análisis Funcional; D: Datos; S: Seguridad; Q: Calidad; V: Desarrollo; I: Arquitectura/Integración; O: Operación/SRE; L: Implantación; J: Proyecto. Los códigos \texttt{TIT} conservan a los responsables nominados en SD1. Los códigos \texttt{REF} identifican puestos por proveer, sin afirmar que exista contratación, acreditación o disponibilidad confirmada. Un mismo código de refuerzo representa un puesto dentro de cada período; su continuidad entre meses se confirmará en la nómina. El responsable de aceptación del paquete conserva la responsabilidad indicada en T-14 aunque varios ejecutores contribuyan a sus actividades.

\paragraph{Capacidad y funciones compartidas.}
Se fija para esta propuesta un máximo de 50 HH de paquetes fijos por persona y quincena, equivalente a 100 HH por mes. Frente a la referencia inicial de 144 HH productivas, las 44 HH restantes deben absorber trabajo variable, observaciones, reensayos y otros compromisos; no constituyen una reserva íntegra disponible ni permiten dar por cubierta una demanda desconocida. La disponibilidad efectiva se reducirá por ausencias y calendario real. Se suman las horas I del titular de Arquitectura/Integración, sin duplicar su jornada. José Lara Arce dispone de 50 HH quincenales conjuntas: O desde el mes 6; entre 8 y 20, hasta 25 HH O y 25 HH L; desde 21, hasta 50 HH O. La carga L posterior se asigna a un puesto de refuerzo, bajo su coordinación. Esta distribución de paquetes no modifica los porcentajes de dedicación comprometidos en SD1.

\paragraph{Regla de asignación y alcance.}
La secuencia de ejecución directa traslada once consumidores de la primera a la segunda quincena de su mismo mes y conserva las demás quincenas. Los ensayos mantienen las ventanas previstas. Se asigna una contribución de perfil a un ejecutor y se comprueba su capacidad diaria; Calidad se ejecuta después de las demás contribuciones del paquete. El registro \path{secuencia_fija_recursos.csv} conserva inicio, fin y duración en días hábiles relativos, y \path{asignacion_fija_recursos.csv} vincula sus 2.680 contribuciones con ejecutores. Se utiliza un máximo diario de 7,2 HH, derivado de 144 HH productivas en veinte días, manteniendo 50 HH por quincena y 100 por mes. La reserva restante se distribuye en los días con disponibilidad; no se presupone disponible durante una jornada ya ocupada. No se fracciona una misma contribución entre varios ejecutores.

\begingroup\small\setlength{\tabcolsep}{2pt}
\begin{longtable}{R{8mm}R{15mm}R{9mm}R{9mm}R{9mm}R{9mm}R{9mm}R{9mm}R{9mm}R{9mm}R{9mm}R{12mm}}
\caption{Puestos de refuerzo para la secuencia fija por quincena}\label{tab:sd7-refuerzo-fijo}\\\hline
Mes & HH fijas & A & D & S & Q & V & I & O & L & J & Total\\\hline\endfirsthead
\hline Mes & HH fijas & A & D & S & Q & V & I & O & L & J & Total\\\hline\endhead
1 & 88 & 0 & 0 & 0 & 0 & 0 & 0 & 0 & 0 & 0 & 0\\\hline
2 & 232 & 0 & 1 & 0 & 0 & 0 & 0 & 0 & 0 & 0 & 1\\\hline
3 & 864 & 3 & 3 & 1 & 0 & 0 & 1 & 0 & 0 & 0 & 8\\\hline
4 & 320 & 0 & 0 & 0 & 0 & 0 & 1 & 1 & 0 & 0 & 2\\\hline
5 & 416 & 0 & 0 & 1 & 1 & 0 & 1 & 2 & 0 & 0 & 5\\\hline
6 & 544 & 0 & 0 & 1 & 1 & 0 & 1 & 2 & 0 & 0 & 5\\\hline
7 & 224 & 0 & 2 & 0 & 1 & 0 & 0 & 0 & 0 & 0 & 3\\\hline
8 & 1904 & 0 & 5 & 1 & 5 & 22 & 2 & 0 & 0 & 0 & 35\\\hline
9 & 1400 & 2 & 3 & 1 & 5 & 12 & 1 & 0 & 0 & 0 & 24\\\hline
10 & 240 & 1 & 2 & 1 & 4 & 0 & 1 & 1 & 0 & 0 & 10\\\hline
11 & 744 & 2 & 2 & 1 & 5 & 1 & 2 & 2 & 2 & 0 & 17\\\hline
12 & 160 & 2 & 1 & 0 & 2 & 0 & 0 & 0 & 1 & 0 & 6\\\hline
13 & 648 & 3 & 2 & 1 & 1 & 0 & 0 & 1 & 0 & 0 & 8\\\hline
14 & 320 & 3 & 1 & 0 & 1 & 0 & 0 & 0 & 0 & 0 & 5\\\hline
15 & 1280 & 2 & 2 & 1 & 5 & 7 & 2 & 1 & 1 & 0 & 21\\\hline
16 & 1192 & 2 & 3 & 1 & 5 & 10 & 0 & 2 & 0 & 0 & 23\\\hline
17 & 1384 & 3 & 2 & 1 & 5 & 7 & 2 & 2 & 2 & 0 & 24\\\hline
18 & 552 & 1 & 0 & 1 & 4 & 0 & 1 & 1 & 1 & 0 & 9\\\hline
19 & 388 & 3 & 2 & 0 & 5 & 0 & 0 & 1 & 1 & 0 & 12\\\hline
20 & 568 & 3 & 2 & 1 & 4 & 0 & 0 & 2 & 2 & 0 & 14\\\hline
21 & 368 & 1 & 1 & 0 & 2 & 0 & 0 & 2 & 1 & 0 & 7\\\hline
22 & 152 & 1 & 1 & 0 & 2 & 0 & 1 & 2 & 0 & 0 & 7\\\hline
23 & 152 & 1 & 1 & 0 & 2 & 0 & 1 & 2 & 0 & 0 & 7\\\hline
24 & 232 & 1 & 1 & 0 & 2 & 0 & 1 & 2 & 1 & 0 & 8\\\hline
25 & 152 & 1 & 1 & 0 & 2 & 0 & 1 & 2 & 0 & 0 & 7\\\hline
26 & 264 & 1 & 1 & 0 & 2 & 0 & 1 & 2 & 0 & 0 & 7\\\hline
27 & 152 & 1 & 1 & 0 & 2 & 0 & 1 & 2 & 0 & 0 & 7\\\hline
28 & 152 & 1 & 1 & 0 & 2 & 0 & 1 & 2 & 0 & 0 & 7\\\hline
29 & 152 & 1 & 1 & 0 & 2 & 0 & 1 & 2 & 0 & 0 & 7\\\hline
30 & 200 & 1 & 1 & 0 & 2 & 0 & 1 & 2 & 1 & 0 & 8\\\hline
31 & 152 & 1 & 1 & 0 & 2 & 0 & 1 & 2 & 0 & 0 & 7\\\hline
32 & 320 & 1 & 1 & 1 & 2 & 0 & 1 & 2 & 0 & 0 & 8\\\hline
33 & 152 & 1 & 1 & 0 & 2 & 0 & 1 & 2 & 0 & 0 & 7\\\hline
34 & 152 & 1 & 1 & 0 & 2 & 0 & 1 & 2 & 0 & 0 & 7\\\hline
35 & 152 & 1 & 1 & 0 & 2 & 0 & 1 & 2 & 0 & 0 & 7\\\hline
36 & 232 & 1 & 1 & 0 & 2 & 0 & 1 & 2 & 1 & 0 & 8\\\hline
37 & 152 & 1 & 1 & 0 & 2 & 0 & 1 & 2 & 0 & 0 & 7\\\hline
38 & 264 & 1 & 1 & 0 & 2 & 0 & 1 & 2 & 0 & 0 & 7\\\hline
39 & 152 & 1 & 1 & 0 & 2 & 0 & 1 & 2 & 0 & 0 & 7\\\hline
40 & 152 & 1 & 1 & 0 & 2 & 0 & 1 & 2 & 0 & 0 & 7\\\hline
41 & 152 & 1 & 1 & 0 & 2 & 0 & 1 & 2 & 0 & 0 & 7\\\hline
42 & 200 & 1 & 1 & 0 & 2 & 0 & 1 & 2 & 1 & 0 & 8\\\hline
43 & 152 & 1 & 1 & 0 & 2 & 0 & 1 & 2 & 0 & 0 & 7\\\hline
44 & 320 & 1 & 1 & 1 & 2 & 0 & 1 & 2 & 0 & 0 & 8\\\hline
45 & 152 & 1 & 1 & 0 & 2 & 0 & 1 & 2 & 0 & 0 & 7\\\hline
46 & 152 & 1 & 1 & 0 & 2 & 0 & 1 & 2 & 0 & 0 & 7\\\hline
47 & 152 & 1 & 1 & 0 & 2 & 0 & 1 & 2 & 0 & 0 & 7\\\hline
48 & 232 & 1 & 1 & 0 & 2 & 0 & 1 & 2 & 1 & 0 & 8\\\hline
49 & 152 & 1 & 1 & 0 & 2 & 0 & 1 & 2 & 0 & 0 & 7\\\hline
50 & 264 & 1 & 1 & 0 & 2 & 0 & 1 & 2 & 0 & 0 & 7\\\hline
51 & 152 & 1 & 1 & 0 & 2 & 0 & 1 & 2 & 0 & 0 & 7\\\hline
52 & 152 & 1 & 1 & 0 & 2 & 0 & 1 & 2 & 0 & 0 & 7\\\hline
53 & 296 & 1 & 2 & 0 & 2 & 0 & 2 & 2 & 0 & 0 & 9\\\hline
54 & 400 & 1 & 1 & 1 & 2 & 0 & 2 & 2 & 1 & 0 & 10\\\hline
55 & 200 & 1 & 1 & 0 & 2 & 0 & 2 & 2 & 0 & 0 & 8\\\hline
56 & 368 & 1 & 1 & 1 & 2 & 0 & 2 & 2 & 0 & 0 & 9\\\hline
\end{longtable}\FuenteTabla{Secuencia propuesta de esfuerzo directo, enlaces resueltos y Calidad posterior; puestos por proveer. Máximo 7,2 HH/persona/día y 50 HH/persona/quincena; falta integrar condiciones externas y calendario real.}\endgroup


El pico de esta secuencia propuesta es de 35 puestos adicionales en el mes 8: 22 de Desarrollo, cinco de Datos, cinco de Calidad, dos de Arquitectura/Integración y uno de Seguridad. Sumados a los ocho titulares, se prevé un techo de 43 puestos para ese mes. Los meses 15, 16 y 17 requieren respectivamente 21, 23 y 24 refuerzos. El cálculo incorpora duraciones y el orden de las revisiones. Es el resultado de una asignación reproducible, sin afirmar optimización ni contratación. La suficiencia integral depende además de puertas, esperas, ambientes, trabajo variable y calendario efectivo. Los frentes E1/E2 conservan horas diferenciadas por paquete.

\paragraph{Provisión y control previo a ejecución.}
La Dirección de Proyecto reservará los puestos del período y registrará nombre, competencias, dedicación, ubicación, modalidad, suplencia y costo en T-8 antes de su incorporación. Cada puesto dispondrá de estación, accesos mínimos y ambiente apropiado; los ocho equipos propios de SD1 no cubren por sí solos los puestos adicionales. Calidad mantendrá independencia respecto de la construcción que revisa. Los reemplazos de titulares conservarán la autorización contractual exigida; el refuerzo de ejecución no supone reemplazo automático. La asignación por paquete se revisará quincenalmente junto con las obligaciones de SD1.

El costo de trabajo conserva las tarifas por perfil del diccionario: $C_{\mathrm{fijo}}=\sum_{k,p}h_{k,p}r_p$. Repartir las mismas HH entre titular y refuerzo no agrega nuevamente el costo del paquete. Cuando una tarifa de refuerzo difiera, el presupuesto sustituirá la tarifa de sus HH y sumará por separado reserva de disponibilidad, incorporación, transferencia adicional, equipamiento, viajes y cobertura externa, sin duplicaciones. No se fija un precio numérico sin las tarifas y condiciones de provisión correspondientes.

Antes de fijar la línea base se incorporarán los 328 HH estacionales y los trabajos variables, se completarán las precedencias de actividades, las puertas y las esperas con duración y ejecutor, y se cruzará la programación con congelamiento, eventos, varadas y marejadas. Si alguna actividad no admite reparto paralelo o la reserva resulta insuficiente, se modificará la secuencia o la dotación conservando los hitos contractuales. El balance quincenal elimina la sobreasignación aritmética del alcance fijo distribuido; la factibilidad integral de plazos y cobertura requiere estas comprobaciones adicionales.
